%======================================================================
\section{Method}\label{sec:method}
%======================================================================

The method keeps LightGCN~\citep{he2020lightgcn} as the scoring model and changes two things.
First, the negatives in the Bayesian personalised ranking loss~\citep{rendle2009bpr} are
weighted by the probability that each sampled unlogged pair is a positive. Second, the score
contains a fixed offset built from an imputed outcome. Both quantities come from two nuisance
models that are fitted once on the log before training: an exposure propensity $\hat p$ and an
imputed outcome $\Yhat$. Algorithm~\ref{alg:qp} summarises the procedure. We call the method
\qp. Apart from the shipped Coat propensities, which are discussed below, nothing in this section
uses uniformly collected data. Those data enter only the selection and evaluation described in
Section~\ref{sec:setup}.

\begin{algorithm}[t]
\caption{\qp: training on the biased log with Q-weighted negatives and an imputation offset}\label{alg:qp}
\begin{algorithmic}[1]
\Require biased log $(O,\,O\odot Y)$; job grid $\mathcal J$ over $f_n$, $\beta$, learning rate,
weight decay, width and depth; splits $k=1,\dots,K$ with validation parts $\mathcal V_k$ and test
parts $\mathcal T_k$
\Ensure for each split $k$, the test nDCG of the job, configuration and epoch with the best
validation nDCG on $\mathcal V_k$
\State $\hat p \gets$ logistic exposure model fitted to $O$ (on Coat, the shipped propensities),
truncated below at $0.05$ \Comment{log only}
\State $\Yhat \gets$ additive model fitted to $O\odot Y$ with inverse-propensity weights $O/\hat p$,
$\lambda=5$, 20 sweeps, clipped to $[0,1]$ \Comment{log only}
\State $Q_{ui} \gets \Yhat_{ui}(1-\hat p_{ui})/(1-\hat p_{ui}\Yhat_{ui})$ for all pairs, clipped to $[0,1]$
\State $z_{ui} \gets$ logit of $\Yhat_{ui}$ (clipped to $[10^{-3},1-10^{-3}]$), standardised over all pairs
\For{each job $g\in\mathcal J$}
  \State train LightGCN $f_\theta$ with the loss of Eq.~\eqref{eq:loss} and score $f_\theta+\beta_g z$
  on the logged positives
  \State after each epoch, compute nDCG on the validation and test parts of every split
  (test values are stored and never used for selection)
\EndFor
\For{each split $k$}
  \State select the job, configuration and epoch with the best validation nDCG on $\mathcal V_k$
  \State report the test nDCG of the selected model on $\mathcal T_k$
\EndFor
\end{algorithmic}
\end{algorithm}

\subsection{Nuisance models}\label{sec:nuisance}

\paragraph{Exposure propensity.}
On Coat we use the propensities shipped with the dataset. On Yahoo!\,R3, KuaiRand-Pure and
KuaiRec we fit a logistic exposure model on the log,
\begin{equation*}
\hat p_{ui}=\sigma\bigl(a+b\,\tilde\ell_u+c\,\tilde\ell_i\bigr),
\end{equation*}
where $\tilde\ell_u$ and $\tilde\ell_i$ are standardised logits of the user and item exposure
rates (the fraction of items logged for each user, and the fraction of users who logged each
item). The coefficients are fitted by Newton's method on all $mn$ pairs of the log. Every
propensity is truncated below at $0.05$, including the shipped ones. The model uses the
exposure matrix $O$ only.

\paragraph{Imputed outcome.}
We fit the additive model $\Yhat_{ui}=\min\{1,\max\{0,\mu+b_u+b_i\}\}$ to the logged outcomes
$O\odot Y$, weighting each observed pair by $1/\hat p_{ui}$. The parameters are updated in closed
form by 20 alternating sweeps of ridge regression with penalty $\lambda=5$. The fit uses the
log only, and it is not cross-fitted, so the same log supplies $\Yhat$ and the training data.
The truncation level, the penalty and the number of sweeps are fixed in the code. They are not
selected on the validation or test data.

On Coat the shipped propensities were estimated by the dataset authors from the uniform
ratings~\citep{schnabel2016}. These are the same uniformly collected ratings that serve for
validation and testing. Coat is therefore the one dataset where uniform data reach a training
nuisance, indirectly. We report Coat with this caveat, and it is listed among the limitations.
The same shipped propensities enter every propensity-based baseline on Coat (DR-LightGCN, NAVIP, DR-JL, MRDR,
DR adjacency and DRUP), so the Coat comparison shares this input.

\subsection{Posterior of an unlogged pair}\label{sec:posterior}

Under the notation of Section~\ref{sec:setting}, a pair with $O_{ui}Y_{ui}=0$ is one whose outcome is not
revealed by the log. It is either unexposed or exposed with a negative outcome. BPR draws its negatives
from this set. The probability that such a pair is a positive is
\begin{equation}\label{eq:posterior}
Q_{ui}=\Pr\bigl(Y_{ui}=1\mid O_{ui}Y_{ui}=0,\,x_{ui}\bigr)=\frac{\Yhat_{ui}\,(1-\hat p_{ui})}{1-\hat p_{ui}\,\Yhat_{ui}}.
\end{equation}
Here $\hat p$ denotes the truncated propensity $\bar p$ that the code uses. The conditioning in $Q$ is on the
covariates $x_{ui}$ and on $O_{ui}Y_{ui}=0$ only, not on $O_{ui}$. The derivation assumes that, given the covariates $x_{ui}$ on which the propensity depends,
exposure is independent of the outcome, and that $\Yhat_{ui}$ estimates $\Pr(Y_{ui}=1\mid x_{ui})$.
Then $\Pr(Y_{ui}=1,\,O_{ui}=0)=\Yhat_{ui}(1-p_{ui})$ and $\Pr(O_{ui}Y_{ui}=0)=1-p_{ui}\Yhat_{ui}$.
The closed form therefore does not cover an exposure mechanism that depends on the outcome beyond
those covariates; Section~\ref{sec:setting} states this as Assumption~\ref{as:indep}. We treat the independence as a
modelling assumption and do not test it.

Three properties follow. First, $0\le Q_{ui}\le\Yhat_{ui}$. Second, $Q_{ui}\to\Yhat_{ui}$ as
$\hat p_{ui}\to 0$, so a pair that is almost never exposed keeps its prior. Third,
$Q_{ui}\to 0$ as $\hat p_{ui}\to 1$ when $\Yhat_{ui}<1$, so a pair with an estimated exposure propensity near one
and no logged positive keeps its full weight as a negative. The quantity has the same form as the exposure
posterior of ExpoMF~\citep{liang2016exposure}. Here it is computed in closed form from the two
nuisances and needs no latent-variable model. In the code, $Q$ is clipped to $[0,1]$ and the
denominator is bounded below by $10^{-6}$. $Q$ is used as a weight. We do not claim that the
resulting training loss is unbiased.

\subsection{Q-weighted BPR}\label{sec:weighted}

Training uses only the logged positives. For each positive pair $(u,i)$ one negative item $j$ is
drawn uniformly from all items. If $(u,j)$ is itself a logged positive, $j$ is drawn once more.
For a mini-batch $B$ of such triples the loss is
\begin{equation}\label{eq:loss}
\begin{split}
\mathcal L&=-\,\frac{\sum_{(u,i,j)\in B}w_{uj}\,\log\sigma\bigl(s_{ui}-s_{uj}\bigr)}{\sum_{(u,i,j)\in B}w_{uj}}\\
&\quad+\lambda\,\frac{1}{|B|}\sum_{(u,i,j)\in B}\bigl(\|e_u\|^2+\|e_i\|^2+\|e_j\|^2\bigr),\\
w_{uj}&=\max\bigl\{1-f_n\,Q_{uj},\;10^{-3}\bigr\}.
\end{split}
\end{equation}
where $e_u$, $e_i$ and $e_j$ are the propagated embeddings and $\lambda$ is the weight decay. The
weight acts on the negative only. The loss is normalised by the total weight of the batch, so its
scale does not depend on how many negatives are down-weighted. The factor $f_n$ equals 1 in every
final job except the Coat job without offset, where $f_n\in\{0.5,1\}$ is searched. With $f_n=1$
and $Q_{uj}$ close to one, the sampled negative contributes almost nothing.

The score $s_{ui}$ is defined in Section~\ref{sec:offset}. The encoder $f_\theta$ is LightGCN with
symmetric normalisation of the unweighted graph of logged positives and a mean over the layer
outputs. The depth and width are tuned as described in Section~\ref{sec:setup}. Adam is used for
optimisation, and training stops early on each split's validation score.

\subsection{Imputation offset}\label{sec:offset}

Let $z_{ui}=(\ell_{ui}-\bar\ell)/s_\ell$, where $\ell_{ui}=\operatorname{logit}\,\Yhat^{c}_{ui}$,
$\Yhat^{c}$ is $\Yhat$ clipped to $[10^{-3},1-10^{-3}]$, and $\bar\ell$ and $s_\ell$ are the mean
and standard deviation of $\ell$ over all $mn$ pairs. The score is
\begin{equation}\label{eq:offset}
s_{ui}=f_\theta(u,i)+\beta\,z_{ui},\qquad \beta\in\{0,\,0.5,\,1,\,2\}.
\end{equation}
The offset is used in training and at inference. It is added to the sampled negative as well,
so it acts on the margin between a positive and its negative:
$s_{ui}-s_{uj}=f_\theta(u,i)-f_\theta(u,j)+\beta(z_{ui}-z_{uj})$. The offset is fixed, so the
training loss cannot change it. The model therefore learns a residual on top of the imputed
preference. Logit adjustment for long-tailed classification~\citep{menon2021logit} also adds a
fixed prior to the logits, but here the prior is an inverse-propensity estimate of a
population-level preference rather than a class frequency. The value $\beta=0$ removes the offset.
Each split selects $\beta$ on its validation part, as one of the jobs in the grid.

\subsection{Procedure and cost}\label{sec:cost}

Algorithm~\ref{alg:qp} lists the steps. Steps 1 to 4 use the log only. The uniform data are read
only in the evaluation inside the training loop and in the selection that follows, and they never
enter the loss. The propensity, the imputation, $Q$ and $z$ are dense
$m\times n$ arrays, so memory grows with $mn$. The largest catalogue used here is KuaiRec, with
7{,}176 users and 10{,}728 items. Each training step adds one lookup of $Q$ per sampled negative
and one lookup of $z$ per positive and per negative. LightGCN itself is unchanged. For larger
catalogues $Q$ and $z$ could be computed on the fly from the imputation parameters and the
propensity model. We have not implemented this. We do not report run times, because the result
files do not contain a controlled timing comparison.

\subsection{Training-free DR propagation and the walk correction}\label{sec:drup}

DRUP is a training-free walk-corrected propagation operator. In this paper it plays two roles: as a baseline, and as a secondary
component whose correction we test separately. From the log it forms two edge estimates for every
pair. The first is the inverse-propensity estimate $W^{\mathrm{IPS}}=OY/\bar p$. The second is the
doubly robust estimate $W^{\mathrm{DR}}=\Yhat+O(Y-\Yhat)/\bar p$, where $\bar p$ is the truncated
propensity~\citep{dudik2011doubly}. The operator propagates these estimates over three hops. A
walk that reuses an edge contributes a biased product, and DRUP's idempotent correction removes
the excess that such walks add. The score combines the one-hop and the corrected three-hop terms
with a weight that is selected on validation data. Under fixed nuisances, the construction is
designed to make the three-hop operator unbiased for propagation over the fully exposed graph. The
argument is given in the supplement.

The correction does not improve accuracy on real data in our experiments. At matched
configurations, the corrected operator is lower on Coat and close to the uncorrected one on Yahoo!\,R3
(Section~\ref{sec:results}). Its gain appears only in a semi-synthetic setting with true propensities and fixed
degree weights, and in one of the tested settings it reverses at large walk weights. We therefore present the correction
as a secondary, training-free component. The strong accuracy of training-free DR propagation on
KuaiRec is reported as a property of the propagation, not of the correction. \qp\ does not use
DRUP for training.
